\section{About this encyclopaedia}\label{sec:about}

This is a prototype of a living, hypertextual reference work on the topological
dynamics of topological groups, with an emphasis on universal minimal flows. It
is meant to grow: entries are added and revised over time, and every
definition, theorem and remark keeps a permanent \emph{tag}, a four-character
identifier that never changes and is never reused, even when the item moves,
is renumbered or is reworded. Cite results by tag, not by number.

Each statement lists the earlier statements it depends on; together these
dependencies form the dependency graph of the encyclopaedia, which can be
browsed from the navigation menu. The interactive proof tree at
\url{https://topologicaldynamics.org/proof-tree/} draws every statement with
what it rests on; selecting a result shows its statement, its proof and its
proof tree. The proof trees of the main results, the theorems that no other
statement uses, are collected at
\url{https://topologicaldynamics.org/proof-tree/main/}.

\paragraph{Conventions.} Topological groups and compact spaces are Hausdorff.
Actions are left actions. Conventions on left and right uniform structures
differ across the literature; the ones used here are fixed in
Definition~\ref{def:ruc}.

\paragraph{Sources and contributing.} The encyclopaedia is built from the
public repository
\url{https://github.com/eulogica/topdyn-encyclopaedia}. Every section links
to its source files there, with their history. Corrections and new entries are
proposed by pull request and reviewed by the editor; the repository's
\texttt{CONTRIBUTING.md} file explains how.

\paragraph{Status of the seed entries.} The first entries were drafted with
the help of an AI assistant from the cited literature. Statements and
attributions were checked against the original papers' abstracts and
statements and against standard secondary sources, and every bibliographic
record was checked against the publisher or arXiv. They have not yet been
reviewed by the editor or by experts. Sections~\ref{sec:kpt-umf}
to~\ref{sec:ea-catalogue} were added later in the same way, from the full
texts of the cited papers where available and otherwise from their abstracts,
and have not been reviewed either. Sections~\ref{sec:maps-transformations}
to~\ref{sec:inductive-limits} report results from preprints
posted on Zenodo, which have not been refereed. Sections~\ref{sec:proofs-maps-transformations}
to~\ref{sec:proofs-inductive-limits} reproduce the proofs from those preprints,
in the versions cited in the bibliography; the bibliographic data of the works
they cite were checked against Crossref, zbMATH or arXiv.

\paragraph{The editor.} The encyclopaedia is edited by LogicArtificialis
(\url{https://logicartificialis.org}). Its seal and one of its emblems,
reproduced from its website, are shown below.

\begin{center}
\emblem{seal}{A gold circular seal with a sapphire turtle, a corgi, a capybara
  and the Latin motto Organum rationis, sapphirus mentis, above the wordmark
  LogicArtificialis.}\hspace{2em}%
\emblem{meaning-guides}{A corgi and a capybara riding a green turtle above the
  wordmark LogicArtificialis and the motto Meaning guides. Syntax decides.}
\end{center}

\paragraph{Funding.} This work is supported by ERC grant DAT (Definable
Algebraic Topology), 101077154.

\begin{center}
\emblem[2.9cm]{dat-logo}{The DAT logo: the words Definable Algebraic Topology,
  with large magenta initials D, A and T.}\hspace{2em}%
\emblem[2.9cm]{erc-eu-logo}{Funded by the European Union, with the flag of the
  European Union, beside the logo of the European Research Council, established
  by the European Commission.}
\end{center}

Funded by the European Union. Views and opinions expressed are however those of
the author(s) only and do not necessarily reflect those of the European Union
or the European Research Council Executive Agency. Neither the European Union
nor the granting authority can be held responsible for them.
